%! Projections onto Subspaces — Unit 4, Lesson 4.2 — Guided Notes (Answer Key)
\documentclass[10pt]{article}
\usepackage{linalg-article}
\usepackage{linalg-key}   % pulls in -boxes; do NOT also load -boxes
\usetikzlibrary{calc}
\newcommand{\TallMath}[1]{$\displaystyle #1\rule[-1.4em]{0pt}{3.2em}$}
\newcommand{\vv}{\mathbf{v}}
\newcommand{\ww}{\mathbf{w}}
\newcommand{\xx}{\mathbf{x}}
\newcommand{\bb}{\mathbf{b}}
\newcommand{\pp}{\mathbf{p}}
\newcommand{\ee}{\mathbf{e}}
\newcommand{\avec}{\mathbf{a}}
\newcommand{\zero}{\mathbf{0}}
\newcommand{\T}{\mathsf{T}}

\begin{document}

\pageheader{Unit 4, Lesson 4.2}{Guided Notes --- Answer Key}
\namedateperiod

\vspace{0.10in}

\begin{objectivebox}
\textbf{By the end of this lesson, I will be able to\dots}
\begin{itemize}\setlength{\itemsep}{2pt}
  \item find the projection of a vector onto a \emph{line} and locate the perpendicular error;
  \item project onto a \emph{subspace} by solving the normal equations $A^{\T}A\hat{\xx}=A^{\T}\bb$;
  \item explain \emph{why} the closest point is the one whose error is perpendicular to the space.
\end{itemize}
\end{objectivebox}

\vspace{0.04in}

\begin{vocabbox}
\small
Fill in each term as we build it together.
\vspace{2pt}
\vocabans{Projection $\pp$ (closest point)}{the point \emph{in} the line or subspace that is closest to $\bb$; $\pp=\hat{x}\avec$ for a line, $\pp=A\hat{\xx}$ for a subspace}
\vocabans{Error $\ee=\bb-\pp$}{what is left over after projecting; it is \emph{perpendicular} to the whole space ($\avec\cdot\ee=0$, or $A^{\T}\ee=\zero$)}
\vocabans{Normal equations}{$A^{\T}A\hat{\xx}=A^{\T}\bb$: solve for the weights $\hat{\xx}$, then the projection is $\pp=A\hat{\xx}$}
\end{vocabbox}

\begin{hookbox}
A drone hovers at a point $\bb$ off to the side of a straight road (a line through the origin).
\begin{itemize}\setlength{\itemsep}{1pt}
  \item Which point of the road is \emph{closest} to the drone? \ansline{The foot of the perpendicular --- the drone's ``shadow'' straight down onto the road.}
  \item How do you \emph{know} it is the closest --- what is special about the segment from that point to the drone? \ansline{That segment meets the road at a right angle; any other point is farther (a hypotenuse instead of a leg).}
\end{itemize}
\end{hookbox}

\begin{notesbox}{1. Closest Means Perpendicular Error}
A target $\bb$ sits off a line through the vector $\avec$. The closest point $\pp$ on the line is the one
where the leftover $\ee=\bb-\pp$ points \emph{straight at} the line --- that is, $\ee$ is
\ans{perpendicular} to $\avec$. Every other point on the line is \emph{farther} away (its distance is a
hypotenuse; the perpendicular distance is a leg). We call $\pp$ the \textbf{projection} of $\bb$ and
$\ee=\bb-\pp$ the \textbf{error}.

\begin{center}
\begin{tikzpicture}[scale=1.0, >=stealth]
  % the line through a
  \draw[burgundy, thick] (-0.6,-0.3)--(4.8,2.4) node[right, font=\footnotesize]{line through $\avec$};
  % vectors
  \coordinate (O) at (0,0);
  \coordinate (P) at (3.2,1.6);   % projection p (on line)
  \coordinate (B) at (2.0,3.0);   % target b (off line)
  \draw[->, very thick, royalblue] (O)--(B) node[above left, font=\footnotesize]{$\bb$};
  \draw[->, very thick, burgundy] (O)--(P) node[below right, font=\footnotesize]{$\pp=\hat{x}\avec$};
  \draw[->, thick, charcoal, dashed] (P)--(B) node[midway, right, font=\footnotesize]{$\ee=\bb-\pp$};
  % right-angle mark at P
  \draw[charcoal, thin] ($(P)+(-0.28,0.14)$)--($(P)+(-0.14,0.42)$)--($(P)+(0.14,0.28)$);
  \fill (O) circle (1.4pt) node[below left, font=\footnotesize]{$\zero$};
\end{tikzpicture}
\end{center}
\end{notesbox}

\begin{notesbox}{2. Projection onto a Line}
Since $\pp$ lies on the line, it is some multiple of $\avec$: write $\pp=\hat{x}\,\avec$ for an unknown
number $\hat{x}$. Now force the right angle --- $\avec$ is perpendicular to the error:
\[
  \avec\cdot(\bb-\hat{x}\avec)=0 \quad\Longrightarrow\quad \avec\cdot\bb-\hat{x}\,(\avec\cdot\avec)=0.
\]
Solve the one equation for the one unknown:
\[
  \hat{x}=\frac{\avec\cdot\bb}{\avec\cdot{\color{keyred}\avec}}, \qquad
  \pp=\hat{x}\,\avec=\frac{\avec\cdot\bb}{\avec\cdot\avec}\,\avec.
\]

\textbf{Worked example.} Line through $\avec=\begin{bsmallmatrix} 1\\2 \end{bsmallmatrix}$, target
$\bb=\begin{bsmallmatrix} 4\\3 \end{bsmallmatrix}$. From the Warm-Up, $\avec\cdot\bb=10$ and
$\avec\cdot\avec=5$, so
\[
  \hat{x}=\frac{10}{5}={\color{keyred}\mathbf{2}}, \qquad
  \pp=2\begin{bsmallmatrix} 1\\2 \end{bsmallmatrix}=\begin{bsmallmatrix} {\color{keyred}\mathbf{2}}\\[2pt] {\color{keyred}\mathbf{4}} \end{bsmallmatrix},\qquad
  \ee=\bb-\pp=\begin{bsmallmatrix} 4\\3 \end{bsmallmatrix}-\begin{bsmallmatrix} 2\\4 \end{bsmallmatrix}=\begin{bsmallmatrix} 2\\-1 \end{bsmallmatrix}.
\]
\textbf{Check the right angle:} $\avec\cdot\ee=(1)(2)+(2)(-1)={\color{keyred}\mathbf{0}}$. \checkmark
\end{notesbox}

\begin{notesbox}{3. The Projection Matrix (a Shortcut)}
Because $\pp$ depends on $\bb$ in a linear way, there is a single \textbf{projection matrix} $P$ with
$\pp=P\bb$. For a line,
\[
  P=\frac{\avec\,\avec^{\T}}{\avec^{\T}\avec}.
\]
For $\avec=\begin{bsmallmatrix} 1\\2 \end{bsmallmatrix}$: $\avec\avec^{\T}=\begin{bsmallmatrix} 1 & 2\\ 2 & 4 \end{bsmallmatrix}$ and $\avec^{\T}\avec=5$, so
$P=\frac{1}{5}\begin{bsmallmatrix} 1 & 2\\ 2 & 4 \end{bsmallmatrix}$. Then
$P\bb=\frac{1}{5}\begin{bsmallmatrix} 1 & 2\\ 2 & 4 \end{bsmallmatrix}\begin{bsmallmatrix} 4\\3 \end{bsmallmatrix}=\frac{1}{5}\begin{bsmallmatrix} 10\\20 \end{bsmallmatrix}=\begin{bsmallmatrix} 2\\4 \end{bsmallmatrix}=\pp.$
\emph{One matrix projects any target.} And projecting again does nothing: $P\pp=\pp$, because $\pp$ is
already \ans{on} the line.
\end{notesbox}

\begin{notesbox}{4. Projection onto a Subspace --- the Normal Equations}
Now the ``road'' is a whole plane: the column space $C(A)$ of a matrix $A$ (its columns span the plane).
We want $\pp=A\hat{\xx}$ closest to $\bb$, so the error $\ee=\bb-A\hat{\xx}$ must be perpendicular to
\emph{every column} of $A$. Stacking those dot-product equations is exactly $A^{\T}\ee=\zero$:
\[
  A^{\T}(\bb-A\hat{\xx})=\zero \quad\Longrightarrow\quad
  \boxed{\,A^{\T}A\,\hat{\xx}=A^{\T}\bb\,}\quad\text{(the \textbf{normal equations}).}
\]
Solve this ordinary system for $\hat{\xx}$; then $\pp=A\hat{\xx}$.

\vspace{3pt}
\textbf{Worked example.} Project $\bb=\begin{bsmallmatrix} 6\\0\\0 \end{bsmallmatrix}$ onto the plane
$C(A)$ for $A=\begin{bsmallmatrix} 1 & 0\\ 1 & 1\\ 1 & 2 \end{bsmallmatrix}$.
\[
  A^{\T}A=\begin{bsmallmatrix} 3 & 3\\ 3 & 5 \end{bsmallmatrix},\qquad
  A^{\T}\bb=\begin{bsmallmatrix} 6\\ 0 \end{bsmallmatrix}\ \Rightarrow\
  \begin{bsmallmatrix} 3 & 3\\ 3 & 5 \end{bsmallmatrix}\hat{\xx}=\begin{bsmallmatrix} 6\\ 0 \end{bsmallmatrix}
  \ \Rightarrow\ \hat{\xx}=\begin{bsmallmatrix} {\color{keyred}\mathbf{5}}\\[2pt] {\color{keyred}\mathbf{-3}} \end{bsmallmatrix}.
\]
Then $\pp=A\hat{\xx}=\begin{bsmallmatrix} 5\\ 2\\ -1 \end{bsmallmatrix}$ and
$\ee=\bb-\pp=\begin{bsmallmatrix} 1\\ -2\\ 1 \end{bsmallmatrix}$. The error is perpendicular to both
columns: $\ee\cdot(1,1,1)=0$ and $\ee\cdot(0,1,2)=0$. So $\ee$ lands in the \textbf{left nullspace}
$N(A^{\T})$ --- the \emph{orthogonal complement} of the column space (Lesson~4.1). Projection splits
$\bb$ into a piece \emph{in} the space ($\pp$) plus a piece \emph{perpendicular} to it ($\ee$).
\end{notesbox}

\begin{practicebox}
Show your work.
\begin{enumerate}[leftmargin=1.6em, itemsep=6pt]
  \item Project $\bb=\begin{bsmallmatrix} 1\\2\\3 \end{bsmallmatrix}$ onto the line through
        $\avec=\begin{bsmallmatrix} 1\\1\\1 \end{bsmallmatrix}$. Find $\hat{x}$, $\pp$, and $\ee$, and
        check $\avec\cdot\ee=0$.
        \ansline{$\avec\cdot\bb=1+2+3=6$, $\avec\cdot\avec=3$, so $\hat{x}=2$; $\pp=(2,2,2)$; $\ee=\bb-\pp=(-1,0,1)$; $\avec\cdot\ee=-1+0+1=0$. \checkmark}
  \item In a projection, what does the error $\ee$ represent, and why must it be perpendicular to the
        space?
        \ansline{$\ee$ is the shortest gap between $\bb$ and the space --- the leftover $\bb-\pp$. If it were \emph{not} perpendicular, sliding $\pp$ along the space would shorten the gap, so $\pp$ would not be closest. Perpendicular $\Leftrightarrow$ shortest.}
\end{enumerate}
\end{practicebox}

\begin{teachernote}
The single idea is ``closest $=$ perpendicular error,'' turned into one equation: $\avec\cdot\ee=0$ for a
line (one unknown), $A^{\T}\ee=\zero$ for a subspace (one equation per column $\Rightarrow$ the normal
equations). Common slips: flipping the ratio to $\frac{\avec\cdot\avec}{\avec\cdot\bb}$, and trying to
solve $A\hat{\xx}=\bb$ directly (it usually has \emph{no} solution --- that is exactly why we project).
Section~4's matrix is the least-squares matrix students will meet again in 4.3, where the same $\pp$
becomes the best-fit line; here we treat it purely as ``project $\bb$ onto a plane.'' The error
$\ee=(1,-2,1)$ living in $N(A^{\T})$ is the direct payoff of Lesson~4.1.
\end{teachernote}

\end{document}
